
\documentclass{article}
\usepackage{colortbl}
\usepackage{makecell}
\usepackage{multirow}
\usepackage{supertabular}

\begin{document}

\newcounter{utterance}

\centering \large Interaction Transcript for game `cladder', experiment `full\_v1.5\_default', episode 2504 with Qwen3.5{-}4B.
\vspace{24pt}

{ \footnotesize  \setcounter{utterance}{1}
\setlength{\tabcolsep}{0pt}
\begin{supertabular}{c@{$\;$}|p{.15\linewidth}@{}p{.15\linewidth}p{.15\linewidth}p{.15\linewidth}p{.15\linewidth}p{.15\linewidth}}
    \# & \multicolumn{2}{c}{Player} && \multicolumn{2}{c}{Game Master} \\
    \hline

    \theutterance \stepcounter{utterance}  
    & & & \multicolumn{4}{p{0.6\linewidth}}{
        \cellcolor[rgb]{0.9,0.9,0.9}{
            \makecell[{{p{\linewidth}}}]{
                \texttt{\tiny{[P1$\langle$GM]}}
                \texttt{You are an expert in causal inference. The following question is not a typical commonsense query, but rather a meticulously designed question created by a professor specializing in causal inference, intended to assess the students' mastery of the course content.} \\
\texttt{Imagine a self{-}contained, hypothetical world with only the following conditions, and without any unmentioned factors or causal relationships: Kidney stone size has a direct effect on treatment and recovery. Treatment has a direct effect on recovery. For patients who have small kidney stones and not receiving treatment, the probability of recovery is 88\%. For patients who have small kidney stones and receiving treatment, the probability of recovery is 72\%. For patients who have large kidney stones and not receiving treatment, the probability of recovery is 16\%. For patients who have large kidney stones and receiving treatment, the probability of recovery is 4\%. The overall probability of large kidney stone is 41\%. For patients receiving treatment, would it be less likely to see recovery if the patient had received no treatment?} \\
\texttt{Start your answer with 'Yes' or 'No', followed by additional reasoning or evidence to support your explanation.} \\
            }
        }
    }
     \\ \\

    \theutterance \stepcounter{utterance}  
    & & \multicolumn{4}{p{0.6\linewidth}}{
        \cellcolor[rgb]{0.95,0.95,0.95}{
            \makecell[{{p{\linewidth}}}]{
                \texttt{\tiny{[P1$\leftrightarrow$P]}}
                \texttt{{[}DRAFT{]}} \\
\texttt{Yes} \\
\\ 
\texttt{**Reasoning:**} \\
\\ 
\texttt{To determine whether recovery is less likely if a patient had received no treatment compared to receiving treatment, we must compare the **counterfactual probabilities** of recovery for the same individual under two different interventions: \$P(\textbackslash{}text\{Recovery\} | \textbackslash{}text\{Treatment\})\$ and \$P(\textbackslash{}text\{Recovery\} | \textbackslash{}text\{No Treatment\})\$.} \\
\\ 
\texttt{The question asks about the effect of the treatment itself, which implies we are looking at the **Average Treatment Effect (ATE)** or the difference in outcomes caused by the intervention, holding the patient's characteristics (specifically kidney stone size) constant. The overall population statistics (41\% large stones) are irrelevant to this specific causal question because the likelihood of recovery depends entirely on the specific patient's stone size and the treatment they receive.} \\
\\ 
\texttt{Let us analyze the two scenarios provided in the prompt:} \\
\\ 
\texttt{1.  **Scenario A: Patient has small kidney stones.**} \\
\texttt{    *   If treated: Probability of recovery = **72\%**.} \\
\texttt{    *   If not treated: Probability of recovery = **88\%**.} \\
\texttt{    *   *Comparison:* \$88\textbackslash{}\% > 72\textbackslash{}\%\$. In this case, recovery is **more** likely if the patient had received no treatment.} \\
\\ 
\texttt{2.  **Scenario B: Patient has large kidney stones.**} \\
\texttt{    *   If treated: Probability of recovery = **4\%**.} \\
\texttt{    *   If not treated: Probability of recovery = **16\%**.} \\
\texttt{    *   *Comparison:* \$16\textbackslash{}\% > 4\textbackslash{}\%\$. In this case, recovery is **more** likely if the patient had received no treatment.} \\
\\ 
\texttt{**Conclusion:**} \\
\texttt{In both sub{-}populations (small stones and large stones), the probability of recovery is higher when the patient does *not* receive treatment compared to when they *do* receive treatment. Therefore, for any given patient in this hypothetical world, it would be **less likely** to see recovery if the patient had received treatment compared to if they had received no treatment.} \\
\\ 
\texttt{The answer to the question "For patients receiving treatment, would it be less likely to see recovery if the patient had received no treatment?" is **Yes**.} \\
            }
        }
    }
    & \\ \\

    \theutterance \stepcounter{utterance}  
    & & \multicolumn{4}{p{0.6\linewidth}}{
        \cellcolor[rgb]{0.95,0.95,0.95}{
            \makecell[{{p{\linewidth}}}]{
                \texttt{\tiny{[P1$\leftrightarrow$P]}}
                \texttt{{[}SELF{-}FEEDBACK{]}} \\
\texttt{FORMAT: ok} \\
\texttt{RULES: ok} \\
\texttt{STRATEGY: problem} \\
\texttt{REASON: The answer incorrectly concludes "Yes" because it fails to recognize that the question asks if recovery is less likely under treatment *compared to* no treatment, whereas the data shows recovery is actually higher without treatment for both stone sizes.} \\
            }
        }
    }
    & \\ \\

    \theutterance \stepcounter{utterance}  
    & \multicolumn{4}{p{0.6\linewidth}}{
        \cellcolor[rgb]{0.561,0.737,0.561}{
            \makecell[{{p{\linewidth}}}]{
                \texttt{\tiny{[P1$\rangle$GM]}}
                \texttt{No} \\
\\ 
\texttt{**Reasoning:**} \\
\\ 
\texttt{The initial reasoning correctly identified that for both small and large kidney stones, the probability of recovery is higher without treatment (88\% vs 72\% for small; 16\% vs 4\% for large). However, the conclusion was flipped.} \\
\\ 
\texttt{The question asks: "For patients receiving treatment, would it be less likely to see recovery if the patient had received no treatment?"} \\
\texttt{This is asking for the comparison: \$P(\textbackslash{}text\{Recovery\} | \textbackslash{}text\{No Treatment\}) < P(\textbackslash{}text\{Recovery\} | \textbackslash{}text\{Treatment\})\$?} \\
\\ 
\texttt{Let's re{-}evaluate the data:} \\
\texttt{1.  **Small stones:** \$P(\textbackslash{}text\{Recovery\} | \textbackslash{}text\{No Treatment\}) = 88\textbackslash{}\%\$, \$P(\textbackslash{}text\{Recovery\} | \textbackslash{}text\{Treatment\}) = 72\textbackslash{}\%\$. Here, \$88\textbackslash{}\% > 72\textbackslash{}\%\$. Recovery is *more* likely without treatment.} \\
\texttt{2.  **Large stones:** \$P(\textbackslash{}text\{Recovery\} | \textbackslash{}text\{No Treatment\}) = 16\textbackslash{}\%\$, \$P(\textbackslash{}text\{Recovery\} | \textbackslash{}text\{Treatment\}) = 4\textbackslash{}\%\$. Here, \$16\textbackslash{}\% > 4\textbackslash{}\%\$. Recovery is *more* likely without treatment.} \\
\\ 
\texttt{In both cases, the probability of recovery is **higher** if the patient had received no treatment compared to if they received treatment. Therefore, it is **not** less likely to see recovery if the patient had received no treatment; it is actually more likely.} \\
\\ 
\texttt{The correct answer is **No**.} \\
            }
        }
    }
    & & \\ \\

    \theutterance \stepcounter{utterance}  
    & & & \multicolumn{2}{p{0.3\linewidth}}{
        \cellcolor[rgb]{0.9,0.9,0.9}{
            \makecell[{{p{\linewidth}}}]{
                \texttt{\tiny{[GM$|$GM]}}
                \texttt{No} \\
            }
        }
    }
    & & \\ \\

    \theutterance \stepcounter{utterance}  
    & & & \multicolumn{2}{p{0.3\linewidth}}{
        \cellcolor[rgb]{0.9,0.9,0.9}{
            \makecell[{{p{\linewidth}}}]{
                \texttt{\tiny{[GM$|$GM]}}
                \texttt{no} \\
            }
        }
    }
    & & \\ \\

    \theutterance \stepcounter{utterance}  
    & & & \multicolumn{2}{p{0.3\linewidth}}{
        \cellcolor[rgb]{0.9,0.9,0.9}{
            \makecell[{{p{\linewidth}}}]{
                \texttt{\tiny{[GM$|$GM]}}
                \texttt{game\_result = WIN} \\
            }
        }
    }
    & & \\ \\

\end{supertabular}
}

\end{document}
